% ─── Literature Review ───────────────────────────────────────────────────────
% Seven-theme thematic synthesis ending in a three-gap statement.
% Themes: (1) buildings decarbonisation and the policy frame; (2) heat pumps,
% electrification and district heating; (3) hydrogen in buildings; (4) hydrogen
% production, supply, storage and transport; (5) merit order, capacity value and
% missing money; (6) spatial and bottom-up methods; (7) policy, carbon pricing
% and the Switzerland context. Cross-reference to \ref{sec:method} preserved.

This review is organised by theme. It moves from the scale
and policy framing of buildings decarbonisation (Section~\ref{sec:lit:buildings}),
through the two rival heating pathways of electrification (Section~\ref{sec:lit:hp}) and
hydrogen (Section~\ref{sec:lit:h2}) and the hydrogen supply chain that any heating role
would draw on (Section~\ref{sec:lit:supply}), to the merit-order and system-value
literature that this paper extends (Section~\ref{sec:lit:meritorder}), the spatial and
bottom-up methods on which its demand reconstruction rests
(Section~\ref{sec:lit:spatial}), and finally the policy and Switzerland context
(Section~\ref{sec:lit:switzerland}). Each theme closes by feeding the gap statement in
Section~\ref{sec:lit:gap}.

Together, these themes establish the urgency of heating decarbonisation and set out the
technical and economic evidence that hydrogen's competitive role, if any, is severely
constrained. This paper puts that evidence to the test across 29 countries and three merit-order arenas, and the test does not return a uniform confirmation. Once each carrier
is charged its own equipment capital and its own last-mile network, the base-load cost
comparison divides geographically instead of running one way across Europe
(Section~\ref{sec:results}), and the case against a spontaneous hydrogen heating market
rests instead on capital recovery in dispatch and on the asymmetric fiscal treatment of
the two carriers.

Three policy instruments recur throughout the themes below and are defined here
before first use. The extension of the European Union Emissions Trading System to
buildings and road transport (ETS2) prices the fossil incumbent directly from
2027; the REPowerEU framework, the European Union's response to the 2022 energy
crisis, reframed efficiency and electrification as instruments of security of
supply; and the Nomenclature of Territorial Units for Statistics at its third tier
(NUTS3), the sub-national grain at which this paper reconstructs demand, is the
resolution at which heating need and stock heterogeneity become visible. Together
these form the policy and spatial scaffold on which the review and the analysis
rest.

\subsection{Buildings-Sector Decarbonisation: Scale, Drivers and the Policy Frame}
\label{sec:lit:buildings}

The buildings sector accounts for over one-third of European energy-related
emissions and is therefore the decisive test of mid-century climate feasibility;
its decarbonisation must address three interconnected challenges, namely the
structural constraint of slow stock turnover, the heterogeneity of building types
and climates across Europe, and the policy instruments (ETS2, the Energy
Performance of Buildings Directive (EPBD), and REPowerEU) designed to drive the
transition. It is a first-order contributor to global and European energy use and
greenhouse gas emissions, and space heating is its largest single end use
\citep{IPCC2018,Smith2018}. In the European
Union, buildings account for approximately 36\% of energy-related CO$_2$ emissions
\citep{EuropeanCommission2020GreenDeal}, a share large enough that the recent
synthesis literature treats the sector as decisive for whether mid-century climate
targets are attainable at all \citep{Azevedo2022,Calvin2023Ipcc}.

The challenge is structural and not merely technical, which reframes
decarbonisation as a problem of allocation, not of selecting a single
winner. The stock turns over slowly, most of the 2050 building stock already
stands, and embodied as well as operational emissions bind the achievable
trajectory \citep{Rck2019Embodied,Pomponi2020Buildings}. A
planetary-boundaries assessment of European heating options finds that no single
carrier is unambiguously dominant across every environmental dimension, which
frames decarbonisation as a problem of allocating the right technology to the
right context instead of selecting a universal winner \citep{Weidner2023Planetary}.

Decarbonisation strategies for the sector fall into three broad categories that
the literature treats as complementary and not substitutable: demand
reduction through envelope efficiency, fuel switching to low-carbon carriers, and
structural change in district energy systems \citep{Azevedo2022,Mathiesen2019Towards}. The pace required is steep. Comparative scenario
work for Europe shows that the speed of envelope renovation and heating
electrification needed to reach successive climate goals rises sharply with
ambition, and that delay raises the eventual cost \citep{Victoria2022Speed,
Mdua2023Towards}.

The critical barrier to scaling these strategies is the slow rate of deep
renovation, and it is as much economic and behavioural as technical. The European
stock renovates at roughly 1\% per year, far below the pace implied by net-zero
pathways, and deep retrofits remain a small fraction of that flow
\citep{Esser2019renovation,Filippidou2019}. Split
incentives between landlords and tenants, household budget constraints that push
retrofits into partial and staged sequences, and weak information all slow uptake
\citep{Arbulu2022Barriers,Maia2024Household}.
Decision-support and household-choice studies emphasise that the renovation
decision and the heating-technology decision are coupled, because envelope
quality conditions the performance of a heat pump and therefore its economics
\citep{Mel2023Decisionsupport,Lingard2020Residential}.

The European policy frame has tightened around this diagnosis. The European Green
Deal, the 2018 long-term strategy, and the climate-neutrality objective set the
direction \citep{EuropeanCommission2020GreenDeal,EuropeanCommission2018LTS}; the recast EPBD of 2024 then
converted it into binding instruments, including minimum energy performance
standards and a trajectory towards zero-emission buildings
\citep{EPBD2024,Mdua2023Towards}. The 2022 energy crisis and the
REPowerEU response accelerated the agenda, reframing efficiency and
electrification as instruments of security of supply as well as decarbonisation
\citep{REpowerEU2022,EUCommissionSWD2022}.
The official reference projections now carry these policies through to mid-century
\citep{EC2024RefScenario}. The empirical literature on implementation, however,
is cautious. It documents wide divergence in transposition and renovation outcomes
across member states, and historical analyses of EU-28 residential heat confirm
that fuel mixes and consumption levels differ markedly between countries and over
time \citep{Mdua2023Towards,Cyx2020EU28Heat}. This
heterogeneity, recurring across drivers, stock, and policy, is the first reason a
spatially explicit treatment is warranted.

\subsection{Heat Pumps, Electrification and District Heating}
\label{sec:lit:hp}

Heat pumps are the leading near-term pathway for decarbonising building heat, combining
high thermodynamic efficiency with a carbon intensity that falls automatically as the
grid decarbonises \citep{IEA2022heatpumps,HPTEC2026WEO2025}. The reported coefficient of
performance places a well-installed unit several times ahead of any combustion boiler in
delivered-heat terms, and life-cycle comparisons confirm that the advantage holds in
present and projected grids \citep{Famiglietti2022Heat,Dong2023}. Air-source systems are
increasingly cost-competitive with gas boilers in new build and moderate-renovation
contexts, while ground-source systems offer superior seasonal performance at higher
capital cost \citep{Kim2023,Hoseinpoori2022}.

The central qualification in the literature concerns cold-climate and poorly insulated
stock, where the coefficient of performance falls as outdoor temperature drops and the
design heat load rises; field and review evidence nonetheless shows that modern units
retain workable performance at low temperatures, and that the binding constraint is more
often envelope quality and emitter sizing than the device itself
\citep{Gibb2023Coming,Lingard2020Residential}. Whole-system and optimal-control studies
add that smart and hybrid operation can materially raise the value of electrified heat
and ease its peak demand \citep{Hoseinpoori2022,Vivian2019Algorithm}. Adoption, however,
is mediated by policy and consumer preference as much as by economics, and several
studies show that awareness, trust, and the perceived disruption of installation shape
willingness to switch \citep{Kokoni2021Policy,Corbett2023Pumping}.

District heating is the complementary structural pathway, and the literature treats it
as the natural home for waste and ambient heat that individual systems cannot capture
\citep{Mathiesen2019Towards,Lund2018Status}. The fourth-generation paradigm lowers
network temperatures so that large heat pumps, industrial excess heat, and other
low-grade sources can supply the stack economically, displacing combustion at the margin
\citep{Lund2018Status,Mller2019Heat}. The Heat Roadmap Europe (HRE) programme mapped the
spatial economics of distribution and showed that dense urban areas can be served at
competitive cost by networks fed from such sources
\citep{Persson2019Heat,Mller2019Heat}. Cost-allocation and carbon-accounting studies of
cogeneration emphasise that the merit order inside a district system, and the emissions
credited to it, depend sensitively on how heat and power are jointly priced, a point
this paper inherits when it ranks the district-heat stack
\citep{Doroti2021Evaluation,Javanshir2022Abandoning}.

Thermal storage and sector coupling further raise the value of electrified district heat
by allowing it to absorb cheap off-peak power
\citep{Enescu2020Thermal,Vecchi2022Carnot}. Across both the individual and networked
routes, the recurring finding is that electrification, suitably staged and stored, is
the cost-effective backbone of decarbonised heat; the open question the present paper
addresses is what, if anything, is left for a combustible carrier, both at the cold-snap
peak that electrified supply finds hardest to serve and on the annual cost of delivered
heat once each carrier is charged its own capital and its own network.

\subsection{Hydrogen in the Buildings Sector: The Case, the Critique and the Costs}
\label{sec:lit:h2}

Hydrogen's potential role in building heat is among the more contested questions
in the energy-transition literature, and the disagreement runs deeper than terminology. Proponents emphasise infrastructure compatibility, fuel
flexibility and the avoidance of deep electrical reinforcement, arguing that
repurposed gas networks could deliver hydrogen to existing boiler stock with
limited disruption \citep{Dodds2015,Bothe2021}. Technically, hydrogen boilers operate on
principles close to those of natural-gas boilers and can in principle use
retrofitted distribution pipework, subject to material compatibility and safety
limits, while fuel cells offer a higher-efficiency route that supplies heat and
power together where simultaneous demand is high \citep{IEAHydrogenReview2022,Scita2020}. Blending studies
map the intermediate option of admixing hydrogen into the gas grid, but identify
binding technical and regulatory ceilings well below full substitution
\citep{Erdener2022Review,Giehl2023}. Public-perception research
finds that the lower behavioural disruption of a like-for-like boiler swap is part
of the appeal, even as it cautions that acceptance is conditional on cost and
safety assurances \citep{Thomas2023Hydrogen, Gordon2023Divergent}.

Yet a larger and increasingly convergent critical literature disputes this claim on
thermodynamic and cost grounds, and its strands form a single causal chain from which it
concludes that hydrogen is implausible for mass-market space heating. The root cause is
thermodynamic. Hydrogen produced by electrolysis and burned in a boiler delivers only a
small fraction of the useful heat that the same electricity would provide through a heat
pump, so hydrogen heating multiplies the renewable generation and electrolyser capacity
required \citep{Cebon2022,Rosenow2022PipeDream}. The first consequence of that
efficiency multiplier is cost; evidence reviews and independent assessments consistently
show that, across the large majority of analysed cases, hydrogen is more expensive than heat pumps or district heating, and its plausible niche is narrow. A meta-review of 54 independent studies finds none supporting a widespread role for hydrogen in heating \citep{Rosenow2024MetaReview,Rosenow2022PipeDream,Rosenow2020Pathway}.

The International Council on Clean Transportation (ICCT) assessment of EU households in
2050 and the Oxford Institute for Energy Studies (OIES) analysis of the German heating
sector both reach this conclusion, the former across the climates it examines and the
latter about a wholesale hydrogen conversion \citep{ICCT2023h2heating, Dickel2024ET29}.
The second consequence is that the result, while consistent in direction, is brittle in
magnitude. Scenario and sensitivity studies show that hydrogen's apparent role is
acutely sensitive to input assumptions on green-hydrogen cost, carbon pricing, and
infrastructure, so that small parameter changes swing it between marginal and material
\citep{Quarton2019Curious,Marzouk2024Expectations}. The third consequence follows from
the first two; once cost and brittleness are imposed within an explicit system model,
hydrogen is confined to a residual role, as the model-based race between hydrogen and
heat pumps for space and water heating finds, with heat pumps dominating the
cost-effective frontier under central assumptions and hydrogen relegated to residual and
peaking duty \citep{Billerbeck2023Race}.

The dedicated modelling literature on hydrogen in buildings remains limited in
scope and scale, covering mostly single countries or technologies, and it divides
into two strands that share a common boundary: single-country techno-economic
studies, and dispatch- and system-focused work. The first strand, the
single-country techno-economic studies, of Spain \citep{Maestre2021, Maestre2022},
Italy \citep{Masi2024State}, Germany \citep{Huckebrink2022}, Ireland
\citep{Genovese2023} and a sub-Saharan context \citep{Modu2023}, together with
national planning assessments \citep{Jung2023Plan}, each confirm that hydrogen's
economic case depends on green-hydrogen cost trajectories, carbon pricing, and
infrastructure availability, but none extends the comparison to the full European
context. The second strand, the dispatch- and system-focused work, reaches the same
conclusion from the other direction. \citet{Berger2020} conducted a
linear-programming assessment of power-to-gas in cross-sector decarbonisation at
coarse spatial resolution, and \citet{Slorach2021Net} assessed
net-zero heating for the United Kingdom to 2050; both evaluate hydrogen within wider
systems but do not isolate its merit-order position against heat pumps in the
heating stack across diverse climates. Each strand is thus confined to one country,
one technology, or one building, and the two together leave the comparative European
question open.

The analytical gap the present paper addresses is the absence of a study that
prices hydrogen on short-run marginal cost in a full merit-order dispatch across
the entire diversity of European climate and building stock, deriving its heating
share endogenously from where it wins the peak instead of imposing it as a
scenario target. This is the contribution that distinguishes the present analysis
from the literature reviewed above, and it motivates the supply, merit-order and
spatial literatures reviewed next. Because the analysis also charges each carrier the
equipment capital implied by the peak load it must serve, and the last-mile network
its own delivery requires, its base-load comparison departs from the convergent
finding above in a defined group of countries, a divergence taken up in
Section~\ref{sec:results} and traced there to the fiscal treatment of retail
electricity, not to the cost of hydrogen. Figure~\ref{fig:lit_landscape} positions the
comparative studies cited here on these two axes and shows the resulting gap
directly. The cell combining sub-national resolution with a merit-order derived
heating share is empty, and it is the cell this paper fills.

\begin{figure}[H]
    \centering
    \caption{\textbf{The comparative hydrogen-versus-heat-pump literature.} Placed by
    spatial resolution and by whether the hydrogen heating share is imposed or derived.}
    \includegraphics[width=\textwidth]{paper/P_LIT_landscape.pdf}
    \label{fig:lit_landscape}
    \\[2pt]{\footnotesize\textit{Source: Authors' contribution.}}
\end{figure}

\subsection{Production, Supply, Storage and Transport of Hydrogen}
\label{sec:lit:supply}

The production cost of green hydrogen is the foundation of any heating claim,
because any heating role for the carrier draws on its cost and availability, so
the production and supply literature bounds the question. The global
reviews establish the baseline. Electrolytic (green) hydrogen is presently more
expensive than the unabated fossil routes it must displace, but its cost is
projected to fall steeply with electrolyser learning and cheap renewable
electricity \citep{IEAHydrogenReview2022,IEA2025hydrogenreview}. Observatory and
calculator data for Europe place the levelised cost of production and its
sensitivity to electricity price, utilisation, and capital cost on a common
footing \citep{EHO2024hydrogencosts,EHO2024lcohcalculator}.
Techno-economic studies of European and national supply, several using explicit
optimisation or Monte Carlo treatment of uncertainty, broadly agree that green
hydrogen reaches the range of tens of euros per megawatt-hour by mid-century in
favourable locations, while remaining costlier where renewable resources are poor
\citep{Sens2022Cost,Zun2023Cost}. Blue hydrogen from natural gas
with carbon capture and pink hydrogen from nuclear power are characterised as
transitional or niche options whose competitiveness depends on gas prices and
capture performance, and which the European cost literature does not find cheapest
at mid-century \citep{IEAGHG2022blue,Staffell2018Hydrogen}.

The delivered cost of hydrogen is dominated by seasonal storage and transport,
both geographically concentrated. Production cost is only a busbar figure (the cost
at the production gate, before storage and transport). What a heating comparison turns on is the delivered, stored cost at the point of combustion.
Heating demand is acutely seasonal and the cold-snap peak is precisely when a
dispatchable hydrogen unit would run, so the cost of holding hydrogen from summer
production to winter use, and of moving it to where the stock sits, can exceed the
production cost itself; both layers are therefore as spatially variable as
production, and both must be resolved at the same country grain. The
technical-potential mapping of salt caverns shows that low-cost seasonal storage is
geographically concentrated in a minority of European countries, which the present
paper translates into a country-specific storage adder \citep{Caglayan2020caverns}.

Network repurposing and import costs vary by country and by delivery method.
Network studies of the European Hydrogen Backbone and gas-for-climate pathways
quantify the cost of repurposing gas pipelines and building dedicated hydrogen
corridors, and of per-property conversion at the distribution edge
\citep{GasForClimate2020backbone,EHB2023countries}. Trade and import analyses, including the OIES
import-route assessments and the IRENA trade work, characterise the delivered cost
of hydrogen reaching European markets by pipeline and as ammonia, and the
diseconomies of small-scale shipping \citep{OIES2022hydrogenimports,OIES2024hydrogenimportsET32}. Optimisation studies of national supply chains and
seasonal power-to-hydrogen storage show that transport technology choice and the
value of inter-seasonal storage are themselves model outputs sensitive to geography
\citep{Lahnaoui2021Optimization,Petkov2020Powertohydrogen}. The implication
this paper carries forward is that the delivered, stored cost of hydrogen at the
point of combustion, not the busbar production cost, is what a heating or peaking
comparison must use, and that this cost is irreducibly spatial.
The further implication is that hydrogen's heating role can only be assessed
through the merit-order dispatch framework that allocates resources to the unit
running at short-run marginal cost. This framework, reviewed next, is the paper's
analytical core.

\subsection{Merit Order, Capacity Value and the Missing-Money Problem}
\label{sec:lit:meritorder}

The analytical core of this paper is drawn from the electricity-market literature on
short-run marginal cost, capacity value, and missing money, which it applies to heat.
The foundational result is that competitive dispatch prices energy at the short-run
marginal cost of the marginal unit, so that capital, being sunk, does not set the spot
price; a peaking plant that runs few hours must therefore recover its fixed costs from
scarcity rent and not from energy margin \citep{Joskow2008capacity}. When that rent is
insufficient, the plant faces the classic missing-money problem
\citep{Joskow2008capacity} and will not be built without a capacity mechanism
\citep{Competition2021}.

The carbon-pass-through literature \citep{Fabra2009etspassthrough} establishes the
complementary mechanism that this paper finds recurring. An emissions price is passed
into the marginal cost of fossil generation and thereby into the market price, so that
it is the carbon charge on the fossil competitor, more than any fall in the price of the
clean alternative, that reorders the stack. Studies of technology value in decarbonising
power systems show that the system worth of a firm, dispatchable resource lies in the
scarce hours it covers, not in the energy it supplies, which motivates separating a
capacity role from an energy role
\citep{Bistline2020ValueOfTechnology,Brown2018Synergies}.

A direct line of work asks where hydrogen and related power-to-X options sit in this
picture, and its findings frame the present analysis. The same conclusion, that hydrogen
turbines and electrolytic flexibility are valuable chiefly for covering rare
residual-load peaks and bridging low-renewable spells, a capacity rather than an energy
contribution, recurs across model scopes. For highly renewable Europe, the
sector-coupled optimisation of \citet{Brown2018Synergies} and the heat-demand-peak study
of \citet{Zeyen2021Mitigating} both place hydrogen in this residual role; for Great
Britain, the integrated value-chain optimisation of \citet{Quarton2019Value} reaches the
same finding at the national scale; and at the global scale, the fully renewable
transition pathways of \citet{Bogdanov2019Radical} and \citet{Victoria2021Solar}
likewise confine dispatchable hydrogen to bridging the residual peak. That the result
holds from the national through the continental to the global scope is what makes it
robust to local conditions. The complementary peak-mitigation and
climatological-adequacy literature shows that the depth of the cold, low-wind event (the
Dunkelflaute) governs how much firm capacity the system needs, and that this need is
robust but bounded \citep{Zeyen2021Mitigating,Mellot2024Mitigating}.

Heat-specific applications of the same logic, including studies of peat displacement and
thermal storage in district systems, confirm that dispatchable units at the top of a
heat stack run for limited hours and so face the same recovery problem
\citep{Javanshir2022Abandoning,Enescu2020Thermal}. To operationalise this framework, the
paper draws its cost and performance parameters from standard technology databases,
European electricity-market reviews and grid-intensity indicators, and carbon-price
trajectories from ETS2 modelling, with the sources and assumptions detailed in
Section~\ref{sec:method}. The European market reviews and grid-intensity indicators that
anchor present-day electricity prices and carbon intensities are by now a settled
empirical base for such comparisons \citep{EMBER2024electricity, EEA2024gridintensity}.
This paper's contribution is to carry the merit-order, capacity-value and missing-money
triad, established for power, into the heating peak and the district-heat stack, and
then back into the power sector that electrified heat runs on, applying one consistent
operating-cost test across all three arenas.

\subsection{Spatial and Bottom-Up Methods for Building-Stock Modelling}
\label{sec:lit:spatial}

Regional heterogeneity is a first-order determinant of both heating demand and
technology feasibility, so the method by which demand is built matters as much as
the technologies compared. Climate zones, vintage distributions, urbanisation and
national regulation all vary substantially across Europe, and aggregate national
models can obscure the spatial targeting that effective policy requires
\citep{Hotmaps2020,Mller2019Heat}. The
bottom-up tradition reconstructs demand from physical stock characteristics, where the top-down tradition infers it from national totals, and reviews of the approach set out its
strengths and its data demands \citep{SueWing2008Synthesis,Prina2020Classification}. Three data infrastructures underpin the present reconstruction.
The Hotmaps project provides an open NUTS3 heat-demand dataset and toolbox derived
from stock and climate data, which this paper uses as an independent validation
reference \citep{Hotmaps2020,Hotmaps2019}.
The EUBUCCO database assembles harmonised footprints and attributes for more than
200 million individual European buildings, supplying the geometric basis for a
footprint-up demand estimate, and is complemented by emerging global building
datasets \citep{MilojevicDupont2023Eubucco, Zhu2025GlobalBuildingAtlas}. The
TABULA typologies translate stock into energy by attaching per-class, per-cohort
space-heating and hot-water intensities across European countries, with national
brochures providing the country detail \citep{Loga2016Tabula,Cyx2011TabulaBelgium}.

Census, geographic, and observatory sources complete the reconstruction and its
benchmarking. Dwelling counts and building-type splits come from the European census and
national equivalents, with the NUTS reference geographies aligning the spatial units
\citep{Eurostat2021census,ONS2021TS044}. The EU Building Stock Observatory and the
Odyssee-MURE profiles provide independent national totals and intensities for
multi-source reconciliation, alongside the peer-reviewed historical series for EU-28
residential heat \citep{BSO2025}. Heating demand is driven by climate through
heating degree days, whose European datasets and projected changes the present analysis
uses for the climate multiplier and warming sensitivity
\citep{EurostatHDD2025}. Recent work disaggregating energy consumption and
emissions to municipal and sub-national resolution for individual countries demonstrates
both the feasibility and the value of high-resolution stock modelling, while remaining national in coverage \citep{Patil2025Disaggregated}. On the network side, sub-national European hydrogen work does exist and is continental in scope \citep{Neumann2023H2Network,Kountouris2024Unified}; what it returns is a network plan rather than a per-market balance sheet for a single plant.

On the choice and dispatch side, the technology-share step in this paper draws on
discrete-choice theory \citep{McFadden1974ConditionalLogit, Train2009DiscreteChoice},
the least-cost layer on linear-optimisation generators and reviews of bottom-up
energy-system models \citep{Krien2020Oemofsolpha,Juan2021Review}, the multi-criteria
reconciliation on compromise-programming methods \citep{Yu1973GroupDecision,
Zeleny1973CompromiseProgramming}, the demand decomposition on the LMDI index method
\citep{Ang2005}, and the uncertainty treatment on reviews of uncertainty assessment and
time-series aggregation in energy-system optimisation \citep{Yue2018Uncertainty,
Hoffmann2020Review}. What is missing from this otherwise rich methodological literature
is a study that combines a footprint-up demand reconstruction at NUTS3 resolution with
an explicit merit-order test of hydrogen across the full European country set; the
methods exist in parts but have not been assembled for this question.

\subsection{Policy, Carbon Pricing and the Switzerland Context}
\label{sec:lit:switzerland}

The policy and carbon-pricing literature supplies the levers that drive this
paper's scenarios, and the recurring mechanism it documents is that pricing the
fossil incumbent, more than subsidising the alternative, reorders the heating
stack. ETS2 is the central instrument. Modelling of its price trajectory spans a
wide range depending on complementary policies, and the official framework and
independent projections place it at the heart of buildings decarbonisation after
2027 \citep{EuropeanCommission2024ETS2,Gunther2025ets2}. National measures increasingly act on the same margin. Boiler
phase-out schedules, building-energy acts, and future-homes standards remove the
fossil option directly, complementing the price signal
\citep{EHPA2025boilerbans,GER2024GEG}.
The security-of-supply literature shows how the 2022 crisis recast gas demand
reduction and electrification as strategic as well as climate objectives, and how
gas-network access and household exposure vary across Europe
\citep{Goldthau2023Eu,Helwig2023Eu}. The price and consumption series that
calibrate the scenarios are drawn from Eurostat and the national regulators
\citep{EurostatNRG2025,Eurostat2025prices}.

Switzerland occupies a distinctive position that the limited literature treats as a case
apart: a non-EU member deeply integrated with EU energy markets, surrounded by countries
with ambitious hydrogen strategies, and with substantial hydropower and nuclear
generation that could support domestic low-carbon hydrogen
\citep{Gabrielli2020,BfE2023}. The Energy Perspectives 2050+ and the national and
cantonal strategies set the decarbonisation pathway, and the 2024 national hydrogen
strategy and climate legislation frame the hydrogen question specifically
\citep{SFOE2023perspectives,EnDK2022plus}. Country-specific studies of Swiss heat
decarbonisation, building-stock thermal performance and the energy-performance gap
establish that the Swiss stock is comparatively well insulated but heavily reliant on
fossil heating, and that electrification and district heating dominate the
cost-effective options \citep{Narula2018Strategies,Streicher2018Assessment}.

The winter-electricity-deficit literature is particularly relevant, since it identifies
the cold-season adequacy gap that any Swiss firm-capacity or hydrogen-peaking role would
have to address, and the domestic tariff data bound the comparison
\citep{Mellot2024Mitigating, ElCom2025tariffs}. The Swiss case in this literature,
however, is examined largely on its own terms. The lacuna is specific. No prior study
has examined whether Switzerland, positioned at the intersection of EU hydrogen
ambitions and its own net-zero commitment, can justify hydrogen imports for residential
heating given its existing low-carbon grid and competing district-heat expansion. The
question of how imported decarbonised hydrogen, priced and stored, would compete in
Swiss building heat under a net-zero-by-2050 constraint, which motivates Research
Question~2, has not been addressed.

\subsection{Research Gap}
\label{sec:lit:gap}

The foregoing review reveals three gaps that this paper addresses, each emerging
from a theme above. First, the buildings, hydrogen, and spatial literatures
together establish that heating demand and technology feasibility are
irreducibly heterogeneous across Europe, yet no continental,
quantitative assessment of hydrogen's role in EU buildings
decarbonisation has been built on a bottom-up heat-demand reconstruction resolved to
NUTS3, in the updated REPowerEU policy
context. Second, the hydrogen-in-buildings and merit-order literatures show that
hydrogen's apparent role is highly sensitive to how it is priced and to local
conditions, yet no prior study has systematically evaluated the
hydrogen-versus-heat-pump trade-off on a consistent operating-cost basis across the
full diversity of European building stock and climate, deriving hydrogen's share
from where it can win the peak instead of imposing it.

Third, Switzerland's specific strategic position, surrounded by
hydrogen-ambitious EU members, with substantial hydroelectric and nuclear
capacity, has not been modelled in a comparative context. Specifically, whether
imported decarbonised hydrogen can compete with electrification in Swiss
residential heating under net-zero constraints remains unaddressed, and it is this
that motivates Research Question~2. This paper
addresses all three gaps through the integrated modelling framework described in
Section~\ref{sec:method}.
